\section{Related Works}

\paragraph{Evaluation Focusing on Memory.} 
Prior work evaluates LLM memorization primarily through long context understanding and recall oriented benchmarks. Early stress test evaluations such as Needle in a Haystack \footnote{https://www.anthropic.com/news/claude-3-family}
 probe a model’s ability to retrieve salient information embedded within extended contexts. Subsequent benchmarks including LongBench \citep{bai2024longbench}, L-Eval \citep{an2024eval}, RULER \citep{hsieh2024ruler}, and $\infty$-Bench \citep{zhang2024bench} systematize this retrieval based evaluation through question answering, summarization, and synthetic retrieval tasks.
More recent efforts extend long context evaluation to conversational or episodic settings. LoCoMo \citep{maharana2024evaluating}, LongMemEval \citep{wulongmemeval}, MemoryAgentBench \citep{ai2025memorybench}, MemoryBench \citep{ai2025memorybench}, and EvoMem \citep{wei2025evo} assess whether models can retain and recall information introduced in the previous interactions. However, these benchmarks primarily evaluate static memorization through post hoc recall using a single query and do not involve an agentic or interactive environment in which memory must be actively used.
In contrast, \ours focuses on LLM agents equipped with explicit memorization mechanisms and evaluates memory usage in sequential multi session agentic settings. Our evaluation emphasizes whether information acquired during earlier interactions can be persistently stored and correctly utilized to support later task execution, reflecting more realistic long-term agent behavior. 

\vspace{-5pt}
\paragraph{Evaluation Focusing on Agentic Abilities.}
A complementary line of work evaluates LLM agents through interactive execution benchmarks that emphasize model reasoning, action selection, and tool use in dynamic environments. Web-based agent environments such as WebShop~\citep{yao2022webshop},  Mind2Web~\citep{deng2023mind2web}, and Mind2Web 2~\citep{gou2025mind2web} assess an agent’s ability to navigate web interfaces, invoke tools, and execute grounded actions in response to web transitions. Coding environments, such as SWE-bench~\citep{jimenez2023swe}, focus on software engineering tasks that require iterative reasoning and tool-mediated code edits to resolve isolated issues. More recent compositional search benchmarks such as BrowseComp~\citep{wei2025browsecomp} and BrowseComp+~\citep{chen2025browsecomp} evaluate agents’ capacity for deep research. MemoryGym~\citep{pleines2025memory} measures within-episode retention in a partially observable control 2D environment. While these benchmarks provide valuable testbeds for evaluating agent execution and reasoning, they are typically formulated as single-session, independent tasks and do not require persistent memory across episodes. As a result, the role of agent memory is not explicitly evaluated. Recent work \citep{zhong2024memorybank, wei2025evo} feeds agentic tasks from above benchmarks in a streaming manner to enable test-time learning. However, unlike our setting, these evaluations do not enforce explicit dependencies across individual tasks. \ours is the first one designed to assess agent memory using sequential subtasks with causal dependencies across sessions. 

Several recent benchmarks highlight the gap between information recall from long conversation history and agentic deployment, but still most evaluate memory via question answering or tool grounding over a fixed history

Mem2ActBench~\citep{shen2026mem2actbench},  MemTrack~\citep{deshpande2025memtrack}, EMemBench~\citep{li2026emembench}, and AgentLongBench~\citep{fang2026agentlongbench} construct long tool-call traces or enterprise-style workflow timelines and test whether agents can retrieve the correct facts or parameters to answer/complete post-hoc follow-up queries. They focus on retrieval from static reasoning traces rather than interdependent task sequences where distilled skills can influence future execution (e.g., learning from inductive problems in formal reasoning in \ours). AgencyBench \citep{li2026agencybench} and Beyond Task Completion~\citep{akshathala2025beyond} incorporate memory into agent execution, but use simple fixed add-and-retrieve tools, prioritizing overall agent capability over systematic evaluation on memory mechanisms. 
In contrast, \ours enforces \textit{cross-task causal dependence} and evaluates memory through \textit{end-to-end sequential task completion}, measuring 
if agents can absorb experiences, acquire new skills, distill reusable knowledge from the past and eventually apply the new skill and understandings to inform future decisions rather than merely recalling previously seen facts.
% if agents can reuse information, acquire new skills, and distill reusable knowledge from the past to inform future decisions rather than merely recalling previously seen facts.


%Overall, \ours differs by \textbf{enforcing a cross-task causal dependence} and measuring memory by \textbf{end-to-end sequential task completion}, capturing whether agents can reuse information, learn new skills, and distill useful understanding from the past in the service of future decisions rather than only recalling it.

%EMemBench and Memory Gym~\citep{li2026emembench,pleines2025memory} evaluate episodic memory skills (e.g., recall/temporal/spatial reasoning) via trajectory-derived QA or partially observable control, rather than requiring memory to support a chain of causally dependent subtasks. Mem2ActBench, MemTrack, and related log-based evaluations~\citep{shen2026mem2actbench,deshpande2025memtrack} construct long interaction/tool traces (or enterprise timelines) and test whether agents can retrieve the right facts or parameters for follow-up requests; this makes evaluation convenient but still centers on targeted retrieval from a largely fixed context. Broader agent assessment suites/frameworks (e.g., AgencyBench and Beyond Task Completion) include memory as a pillar but do not impose the controlled inter-subtask dependence needed to attribute failures to memory usage~\citep{li2026agencybench,akshathala2025beyond}. Overall, while these concurrent works strengthen the case for memory-aware evaluation, \ours differs by \textbf{enforcing cross-session causal dependence} and measuring memory by \textbf{end-to-end sequential task completion}, capturing whether agents can \textbf{accumulate and reuse} information in service of future decisions rather than only recalling it.
